\usepackage[table]{xcolor}
\usepackage[english,italian]{babel}
\usepackage{amsmath, amssymb, pifont, geometry, newlfont, caption, multirow, tabularx, svg, pdfpages, placeins, colortbl, amsmath, float, subcaption, enumitem, hyperref}
\usepackage{array} % Per le tabelle
\usepackage{graphicx} % Per includere immagini (necessario anche per pgfplots/tikz)
\usepackage{tikz} % Motore grafico
\usepackage{pgf-pie}
\usepackage{pgfplots} % Interfaccia per i grafici

\usepackage{listings}
\renewcommand{\lstlistingname}{Listato}
\usepackage{setspace}
\usepackage{csquotes}
\usepackage[
  backend=biber,
  style=numeric,
  sorting=ynt
]{biblatex}
\addbibresource{bibliography.bib}

\hypersetup{
  bookmarksopen=true,
  pdfauthor=Riccardo Valtorta
}

\lstdefinestyle{myStyle}{
  language=SQL,
  basicstyle=\small\ttfamily,
  keepspaces=true
}

\lstdefinestyle{myStyle2}{
  language=SQL,
  basicstyle=\small\ttfamily,
  keepspaces=true,
  numbers = left,
  numbersep=1pt
}

\definecolor{codegreen}{rgb}{0,0.6,0}
\definecolor{codegray}{rgb}{0.5,0.5,0.5}
\definecolor{codepurple}{rgb}{0.58,0,0.82}
\definecolor{backcolour}{rgb}{0.95,0.95,0.92}

\lstdefinestyle{myStyle3}{
  language=SQL,
  backgroundcolor=\color{backcolour},
  commentstyle=\color{codegreen},
  keywordstyle=\color{magenta},
  numberstyle=\tiny\color{codegray},
  stringstyle=\color{codepurple},
  basicstyle=\ttfamily\footnotesize,
  breakatwhitespace=false,
  breaklines=true,
  captionpos=b,
  keepspaces=true,
  numbers=left,
  numbersep=5pt,
  showspaces=false,
  showstringspaces=false,
  showtabs=false,
  tabsize=2
}

\lstdefinestyle{py}{
  language=Python,
  backgroundcolor=\color{backcolour},
  commentstyle=\color{codegreen},
  keywordstyle=\color{magenta},
  numberstyle=\tiny\color{codegray},
  stringstyle=\color{codepurple},
  basicstyle=\ttfamily\footnotesize,
  breakatwhitespace=false,
  breaklines=true,
  captionpos=b,
  keepspaces=true,
  numbers=left,
  numbersep=5pt,
  showspaces=false,
  showstringspaces=false,
  showtabs=false,
  tabsize=2
}

\let\oldincludegraphics\includegraphics
\renewcommand{\includegraphics}[2][]{%
  \fbox{\oldincludegraphics[#1]{#2}}%
}

\newcommand{\getNumeroProblema}{%
  \refstepcounter{table}\thetable%
}

\pgfplotsset{compat=1.18}
\usetikzlibrary{calc}
\usetikzlibrary{fpu} % Necessario per \pgfmathprintnumber

% --- PREAMBOLO (Da inserire all'inizio del documento, prima di \begin{document}) ---
\usepackage{tikz}
\usepackage{pifont} % Per il simbolo della spunta (dingbat)

% Definisco i colori personalizzati (puoi modificarli qui)
\definecolor{successgreen}{RGB}{40, 167, 69} % Verde stile "successo"
\definecolor{failred}{RGB}{220, 53, 69}      % Rosso stile "errore"

% Definisco i comandi per i simboli
% \found = Problema individuato (Spunta Verde)
\newcommand{\found}{\textcolor{successgreen}{\Large\ding{51}}}
% \notfound = Problema NON individuato (X Rossa)
\newcommand{\notfound}{\textcolor{failred}{\Large\ding{55}}}

% -----------------------------------------------------------------------------------
% Per gli scenari
\usepackage{tcolorbox}
\tcbuselibrary{skins}

\newtcolorbox{scenario}[2][]{
  enhanced,
  arc=3pt,
  outer arc=3pt,
  boxrule=0.5pt,
  colback=gray!4,
  colframe=gray!50,
  fontupper=\small\sffamily,
  top=6pt, bottom=6pt, left=8pt, right=8pt,
  title={\small\sffamily\textbf{#2}},
  colbacktitle=gray!18,
  coltitle=black!80,
  titlerule=0.3pt,
  toptitle=3pt,
  bottomtitle=3pt,
  #1
}

% ------------------------------------------------------------------------------------
% Per colorare i links
\hypersetup{
  colorlinks=true, 
  urlcolor=blue, 
  linkcolor=black
}
